\documentclass[12pt,a4paper]{article}
\usepackage[margin=1in]{geometry}
\usepackage{amsmath,amssymb,amsthm,mathtools}
\usepackage{graphicx}
\usepackage{hyperref}
\usepackage{xcolor}
\usepackage{booktabs}
\usepackage{listings}
\usepackage{longtable}
\usepackage{array}
\usepackage{fancyhdr}

\hypersetup{colorlinks=true,linkcolor=blue,citecolor=blue,urlcolor=blue}

\pagestyle{fancy}
\fancyhf{}
\lhead{\small Al-Hindara Project}
\chead{\small v7.7.2 Audit Draft (Post-\S Z)}
\rhead{\small Abdulla M. N. Al-Jabri}
\lfoot{\small ORCID: 0009-0003-3319-3822}
\cfoot{\small \thepage}
\rfoot{\small DOI: 10.5281/zenodo.23924915}
\renewcommand{\headrulewidth}{0.4pt}
\renewcommand{\footrulewidth}{0.4pt}

\lstset{
 basicstyle=\footnotesize\ttfamily,
 breaklines=true,
 frame=single,
 backgroundcolor=\color{gray!6},
 tabsize=4
}

\title{A Four-Frame Exploratory Framework:\\
Numerical Damping, Geometric Ratios, Algebraic Symmetry-Breaking,\\
and Spectral Domain Analysis\\
\large v7.7.2 Audit Release — Numerical Variational Evidence}
\author{Abdulla M. N. Al-Jabri\\ Independent Researcher, Sana'a, Yemen\\ ORCID: 0009-0003-3319-3822 --- \href{mailto:jabri62018@gmail.com}{jabri62018@gmail.com}}
\date{November 3, 2026}

\newtheorem{proposition}{Proposition}
\newtheorem{lemma}{Lemma}
\newtheorem{remark}{Remark}

\begin{document}
\maketitle
\thispagestyle{empty}

\begin{abstract}
This paper consolidates four exploratory frameworks of the Al-Hindara project, presented with an explicit methodological audit (\S Z) that distinguishes between what has been proven, what has been numerically verified, what has been chosen as a model, and what remains open. 
The frameworks are: (1) a numerical damping function; (2) a geometric ratio with exploratory field observations; (3) an algebraic gauging identity; (4) a spectral problem where the kernel $g$ acts as a formal zero mode.

The v7.7.2 corrective path establishes the analytical domain inclusion for the zero mode. Furthermore, it presents a numerical variational audit using a constrained B-spline Rayleigh--Ritz method. The numerical evidence strongly suggests a positive spectral gap ($\lambda_{1,\mathrm{RR}} \approx 1.7685$), but the identification of this value with the continuous self-adjoint realization remains an open analytical question. This document does not claim a mathematical closure of the spectral problem.
\end{abstract}

\noindent\textbf{Keywords:} Sturm--Liouville theory; spectral domain analysis; Rayleigh--Ritz method; spectral gap; B-splines.

\tableofcontents
\vspace{0.5cm}

% ============================================================
\section{Introduction and Methodology}\label{sec:intro}
% ============================================================

The Al-Hindara project explores mathematical structures across numerical, geometric, and algebraic domains. The methodology of this paper separates claims into four categories:
\begin{enumerate}
\item \textbf{Analytically Proven:} Exact identities and algebraic relations derived from explicit definitions, with proof sketches.
\item \textbf{Numerically Verified:} Reproducible computations (e.g., Rayleigh-Ritz eigenvalues), presented as supporting evidence—not as proof.
\item \textbf{Empirically Observed:} Exploratory data presented without claims of universal physical laws.
\item \textbf{Open:} Questions not yet resolved in the current path.
\end{enumerate}

\textbf{The general status of Frame 4 remains open.}

% ============================================================
\section{The Numerical Damping Framework (Frame 1)}\label{sec:numerical}
% ============================================================

\subsection{Definition and Sign Analysis}
We consider the real function $Z(x; x_p)$ defined on $x \ge 1$:
\begin{equation}\label{eq:Z}
Z(x; x_p) = x^{-5} \ln(x) \sin\left(\frac{2\pi}{x}\right) e^{-x/x_p}.
\end{equation}
For $x_p > 0$, the sign of $Z(x)$ is dictated by $\sin(2\pi/x)$. Thus, $Z(x) < 0$ on $(1, 2)$, $Z(2) = 0$, and $Z(x) > 0$ on $(2, \infty)$. The roots are $x=1$ and $x=2$.

\subsection{The Accumulation Function}
Define $A(x) := \int_1^x Z(t; 21)\, dt$. Since $A'(x) = Z(x)$, $A(x)$ strictly decreases on $(1, 2)$, reaches a global minimum at $x=2$, and strictly increases for all $x > 2$.

% ============================================================
\section{The Scaling Bridge (v2.2)}\label{sec:bridge}
% ============================================================

The bridge is a \emph{dimensional scaling} operation:
\begin{equation}\label{eq:bridge}
Q_{\text{scaled}} = \underbrace{z}_{\text{dimensionless}} \times \underbrace{\Lambda_{\text{scale}}}_{\text{units}}.
\end{equation}
It does not prove any mathematical model or convert $Z(x)$ into a physical field.

% ============================================================
\section{Geometric Ratios (Frame 2)}\label{sec:geometric}
% ============================================================

At \(\theta = 22.5^\circ\):
\begin{equation}\label{eq:tan}
\tan(22.5^\circ) = \sqrt{2} - 1 \approx 0.41421356.
\end{equation}

\begin{table}[h!]
\centering
\caption{Exploratory field measurements of crown ratios.}
\begin{tabular}{cccccc}
\toprule
Tree & \(H\) (m) & \(h\) (m) & \(C=H-h\) (m) & \(R=C/H\) & Deviation from \(\sqrt{2}-1\)\\
\midrule
1 & 7.80 & 4.50 & 3.30 & 0.42308 & \(+0.00887\)\\
2 & 8.50 & 5.00 & 3.50 & 0.41176 & \(-0.00245\)\\
3 & 7.20 & 4.20 & 3.00 & 0.41667 & \(+0.00245\)\\
4 & 9.00 & 5.40 & 3.60 & 0.40000 & \(-0.01421\)\\
5 & 7.50 & 4.40 & 3.10 & 0.41333 & \(-0.00088\)\\
\midrule
\multicolumn{3}{r}{\textbf{Mean \(R\):}} & & \textbf{0.41297} & \\
\bottomrule
\end{tabular}
\end{table}

% ============================================================
\section{Algebraic Symmetry-Breaking (Frame 3)}\label{sec:algebraic}
% ============================================================

Let \(u \in (0, 2\pi)\) and define the kernel:
\begin{equation}\label{eq:g}
g(u) := \left(\frac{u}{2\pi}\right)^{5} \ln\!\left(\frac{2\pi}{u}\right) e^{-2\pi/(x_p u)}.
\end{equation}
Let \(p(u) := g'(u)/g(u)\). Define \(\mathcal{H}_g\):
\begin{equation}\label{eq:Hg}
\mathcal{H}_g[f] := -g(u)\frac{d^2}{du^2} \left(\frac{f(u)}{g(u)}\right).
\end{equation}

\begin{proposition}[Gauging Identity]
For any \(Y \in C^2(0, 2\pi)\), \(\mathcal{H}_g[gY] = -gY''\).
\end{proposition}

Define unweighted \(L^2\) normalized states \(\Psi_n(u) = g(u)Y_n(u)/\sqrt{N_n}\), with \(Y_n(u) = \sin(nu/2)\). The numerical integration at \(x_p=21\) yields:
\begin{equation}\label{eq:R}
R_{13/12} = \frac{|\mathcal{A}_{13}|^2}{|\mathcal{A}_{12}|^2} \approx 0.6035155248.
\end{equation}
This ratio is a numerical fingerprint, not a universal constant.

% ============================================================
\section{The Spectral Problem (Frame 4)}\label{sec:spectral}
% ============================================================

\subsection{The Operator and the Formal Zero Mode}
Define the Schr\"odinger-like operator:
\begin{equation}\label{eq:L}
L := -\partial_u^2 + V_{\text{eff}}(u), \qquad V_{\text{eff}}(u) := \frac{g''(u)}{g(u)}.
\end{equation}

\begin{proposition}[Formal Zero Mode]\label{prop:zero}
The kernel \(g\) satisfies \(Lg = 0\) on the open interval \((0, 2\pi)\) as a formal identity. To be an actual eigenfunction, the domain and boundary conditions must be specified.
\end{proposition}

\subsection{Analytical Domain and the Kernel}
Let \(A = \partial_u - p\) with \(p = g'/g\). The quadratic form is \(q_F[f] = \|Af\|_{L^2}^2\). 

\begin{lemma}[Domain Inclusion]\label{lem:domain}
The function \(g\) belongs to the closed form domain \(D(q_F)\), and consequently \(g \in \mathcal{D}(L_F)\) with \(L_F g = 0\).
\end{lemma}
\begin{proof}[Sketch]
\(g \in L^2\) and \(Ag = 0 \in L^2\). By using smooth two-sided cutoff functions \(\chi_\epsilon\), we can approximate \(g\) in the form norm. The cutoff errors are \(\int |\chi_\epsilon' g|^2 du \sim \epsilon^9\) at \(0^+\) and \(\sim \epsilon\) at \(2\pi^-\). Thus \(g \in D(q_F)\), and by the definition of the form-associated operator, \(g \in \mathcal{D}(L_F)\) with \(L_F g = 0\).
\end{proof}

\begin{lemma}[One-Dimensional Kernel]\label{lem:kernel}
\(\dim \ker L_F = 1\).
\end{lemma}

\subsection{Numerical Variational Audit}
A constrained B-spline Rayleigh--Ritz computation was performed for the quadratic form
\[
Q[f]=\int_{\varepsilon}^{2\pi-\varepsilon} |f'(u)-p(u)f(u)|^2\,du,
\]
subject to the orthogonality constraint \(\int_{\varepsilon}^{2\pi-\varepsilon} f(u)g(u)\,du=0\).

The numerical investigation included piecewise Gauss--Legendre quadrature (8 to 32 local nodes), endpoint truncations from \(10^{-2}\) to \(10^{-5}\), B-spline refinements (10 to 40 internal knots), and a comparison between standard B-splines and the weighted representation \(f=g\psi\).

\begin{table}[h!]
\centering
\caption{Constrained Rayleigh--Ritz eigenvalues for $N=40$, $\varepsilon=10^{-3}$.}
\begin{tabular}{ccc}
\toprule
Eigenvalue & Value & Spectral Gap \\
\midrule
\(\lambda_{1,\mathrm{RR}}\) & \(1.76852476\) & (Baseline) \\
\(\lambda_{2,\mathrm{RR}}\) & \(3.91103755\) & \(\Delta_{\mathrm{RR}} \approx 2.14251279\) \\
\(\lambda_{3,\mathrm{RR}}\) & \(6.55896485\) & \\
\bottomrule
\end{tabular}
\end{table}

The constraint residuals were at machine precision (\(\sim 10^{-18}\)). The computed eigenvalues showed strong numerical stability under all tested refinements.

\begin{remark}[Numerical vs. Analytical]
These results provide numerical evidence for a stable, positive first constrained variational eigenvalue. They do not, by themselves, establish the spectrum of the continuous operator \(L\).
\end{remark}

\subsection{Open Spectral Identification Problem}\label{sec:open}

The constrained Rayleigh--Ritz computation yields a stable first positive variational eigenvalue, \(\lambda_{1,\mathrm{RR}}\approx1.7685248\). The remaining analytical question is whether this value converges to the first positive eigenvalue of a specified self-adjoint realization of \(L\).

In particular, the following questions remain open:
\begin{enumerate}
    \item What is the appropriate self-adjoint domain of \(L\) on \((0,2\pi)\)?
    \item Does the selected realization satisfy \(Lg=0\) with \(g\) belonging to its operator domain?
    \item Is the kernel strictly one-dimensional in the continuous setting?
    \item Does the constrained variational problem characterize the first positive eigenvalue of that realization?
    \item Can the numerical spectral gap be established analytically or through rigorous spectral convergence estimates?
\end{enumerate}

% ============================================================
\section{\S Z — Evidence Classification}\label{sec:audit}
% ============================================================

\begin{longtable}{p{4.5cm}p{6cm}p{3cm}}
\caption{Classification of Frame 4 Claims.}\label{tab:audit}\\
\toprule
Claim & Status & Evidence\\
\midrule
\endfirsthead
\toprule
Claim & Status & Evidence\\
\midrule
\endhead
\bottomrule
\endfoot
\(Lg = 0\) & Proven (Formal) & Algebraic identity \\
\(g \in \mathcal{D}(L_F)\) & Proven & Cutoff approximation in form norm \\
\(\dim \ker L_F = 1\) & Proven & Local uniqueness of ODE solution \\
\(\lambda_0 = 0\) & Proven & Result of domain inclusion \\
Compact resolvent & Open & Hardy inequalities failed \\
Spectral gap (Continuous) & Open & Requires domain analysis \\
\(\lambda_{1,\mathrm{RR}}\) (Discrete) & Numerically Verified & Rayleigh-Ritz audit \\
\(R_{13/12} \approx 0.6035\) & Numerically Verified & Quadrature integration \\
\bottomrule
\end{longtable}

% ============================================================
\appendix
\section{Reproducibility — Python Code}\label{app:code}
% ============================================================

The following Python code reproduces the B-spline Rayleigh-Ritz numerical audit. (Note: See the \texttt{all-frames2.ipynb} notebook for the complete implementation).

\begin{lstlisting}[language=Python]
import numpy as np
from scipy.interpolate import BSpline
from scipy.linalg import eigh, qr

PI2 = 2 * np.pi
XP = 21.0

def g_vector(u, xp=XP):
    u = np.asarray(u, dtype=float)
    result = np.zeros_like(u)
    mask = (u > 0) & (u < PI2)
    um = u[mask]
    result[mask] = (um / PI2)**5 * np.log(PI2 / um) * np.exp(-PI2 / (xp * um))
    return result

def p_vector(u, xp=XP):
    u = np.asarray(u, dtype=float)
    result = np.zeros_like(u)
    mask = (u > 0) & (u < PI2)
    um = u[mask]
    result[mask] = 5.0 / um - 1.0 / (um * np.log(PI2 / um)) + PI2 / (xp * um**2)
    return result

# [See the ipynb for the full matrix assembly, piecewise Gauss integration, 
# and constrained eigenvalue solver using QR decomposition]
\end{lstlisting}

% ============================================================
% Footer
% ============================================================
\vfill
\begin{center}
\rule{\linewidth}{0.5pt}
\vspace{0.2cm}
\textbf{Al-Hindara Project — Consolidated Draft v7.7.2 (Post-\S Z Audit)}\\
Author: Abdulla M. N. Al-Jabri\\
ORCID: 0009-0003-3319-3822\\
Release status: v7.0.0 (Audit draft; v7.7.2 corrective path)\\
Archived Publication: Zenodo Record \href{https://zenodo.org/records/23924915}{23924915} | DOI: \href{https://doi.org/10.5281/zenodo.23924915}{10.5281/zenodo.23924915}\\
License: CC BY 4.0 --- \copyright\ 2026 Abdulla M. N. Al-Jabri\\
\vspace{0.2cm}
\rule{\linewidth}{0.5pt}
\end{center}

\end{document}